% Part: sets-functions-relations
% Chapter: size-of-sets-complete

\documentclass[../../../include/open-logic-chapter]{subfiles}

\begin{document}

\olchapter{sfr}{siz}{ગણોનું કદ}

\begin{editorial}
આ પ્રકરણ પરિગણના, ગણનીયતા અને અગણનીયતાની ચર્ચા કરે છે.
કેટલાક વિભાગો બે સ્વરૂપે આપેલા છે: વધુ પ્રાથમિક સ્વરૂપમાં
પરિગણનાને યાદી અથવા $\PosInt$થી આવતું વ્યાપ્ત વિધેય
ગણવામાં આવે છે; વધુ અમૂર્ત સ્વરૂપમાં પરિગણનાને
$\Nat$ સાથેનું એક-એક વ્યાપ્ત વિધેય ગણી વ્યાખ્યાયિત કરાય છે.
\end{editorial}

\olimport{introduction}

\olimport{enumerability}

\olimport{zig-zag}

\olimport{pairing}

\olimport{pairing-alt}

\olimport{non-enumerability}

\olimport{reduction}

\olimport{equinumerous-sets}

\olimport{comparing-size}

\olimport{schroder-bernstein}

\begin{editorial}
નીચેના \olref[sfr][siz][enm-alt]{sec},
\olref[sfr][siz][nen-alt]{sec} અને \olref[sfr][siz][red-alt]{sec}
એ અનુક્રમે \olref[sfr][siz][enm]{sec},
\olref[sfr][siz][nen]{sec} અને \olref[sfr][siz][red]{sec}નાં
વૈકલ્પિક સ્વરૂપો છે. ટિમ બટને પોતાના ઓપન સેટ થિયરી
ગ્રંથમાં ઉપયોગ માટે તે તૈયાર કર્યાં છે. તે થોડાં વધુ
ઉન્નત છે અને ગણસિદ્ધાંતના સંદર્ભમાં વધુ અનુકૂળ એવી
ગણનીયતાની જુદી વ્યાખ્યા વાપરે છે: યાદીમાં મૂકી શકાય
અથવા $\PosInt$થી આવતા વ્યાપ્ત વિધેયનો વિસ્તાર હોય
તેને બદલે, $\Nat$ અથવા તેના આરંભખંડ સાથે એક-એક
વ્યાપ્ત વિધેય હોય એવી વ્યાખ્યા.
\end{editorial}

\olimport{enumerability-alt}
\olimport{non-enumerability-alt}
\olimport{reduction-alt}

\OLEndChapterHook

\end{document}
